\chapter{सदिश}\label{ch:g10:vectors}

\omterm{def:g10:vectors:vector}{सदिश} एक विस्थापन को समेटता है: एक दिशा और एक लंबाई, आरंभिक बिंदु चाहे जो हो।
\omterm{def:g10:vectors:vector}{सदिश} ज्यामितीय तथ्यों की स्वच्छ उपपत्तियाँ देते हैं और, जैसे ही \omterm{def:g10:coordgeom:system}{निर्देशांक}
चित्र में आते हैं, उन्हें छोटी-छोटी गणनाओं में बदल देते हैं। \cref{ch:g10:lines} में और आगे
चलकर \cref{ch:g11:scal} के अदिश गुणनफल के लिए वे एक प्रमुख औज़ार होंगे।

\section{स्थानांतरण और सदिश}

\begin{definition}[सदिश]\label{def:g10:vectors:vector}
दो बिंदु $A$ और $B$ दिए होने पर \emph{सदिश}\index{सदिश} $\vect{AB}$ $A$ से $B$ तक
के विस्थापन को दर्शाता है: उसकी एक \emph{दिशा} होती है (रेखा $(AB)$ की दिशा,
और उस पर चलने का ढंग, $A$ से $B$ की ओर) और एक \emph{लंबाई} (दूरी $AB$,
जिसे $\norm{\vect{AB}}$ भी लिखते हैं)। दो सदिश \emph{बराबर} तब होते हैं जब वे एक ही
विस्थापन समेटते हों: $\vect{AB} = \vect{CD}$ का अर्थ है कि दोनों विस्थापनों को किसी भी बिंदु
पर लगाने से एक ही परिणाम मिलता है।
\end{definition}

\begin{proposition}[बराबर सदिश और समांतर चतुर्भुज]\label{prop:g10:vectors:parallelogram}
$\vect{AB} = \vect{CD}$ ठीक तब है जब $ABDC$ (इसी क्रम में!) समांतर चतुर्भुज हो, चाहे वह चपटा ही
क्यों न हो --- अर्थात् तब जब $[AD]$ और $[BC]$ का \omterm{prop:g10:coordgeom:midpoint}{मध्यबिंदु} एक ही हो।
\end{proposition}
\admitted

\begin{omfigure}
\begin{tikzpicture}[scale=1.0]
  \coordinate (A) at (0,0);
  \coordinate (B) at (2.2,0.8);
  \coordinate (C) at (1.0,-1.2);
  \coordinate (D) at (3.2,-0.4);
  \draw[omExo, thin] (A) -- (C);
  \draw[omExo, thin] (B) -- (D);
  \draw[-Stealth, omDef, thick] (A) -- (B)
    node[midway, above, font=\small] {$\vect{AB}$};
  \draw[-Stealth, omDef, thick] (C) -- (D)
    node[midway, below, font=\small] {$\vect{CD}$};
  \fill (A) circle (1.5pt) node[left, font=\small] {$A$};
  \fill (B) circle (1.5pt) node[right, font=\small] {$B$};
  \fill (C) circle (1.5pt) node[left, font=\small] {$C$};
  \fill (D) circle (1.5pt) node[right, font=\small] {$D$};
\end{tikzpicture}

{\small बराबर \omterm{def:g10:vectors:vector}{सदिश} $\vect{AB} = \vect{CD}$: वही दिशा, वही ढंग, वही लंबाई। चतुर्भुज $ABDC$
समांतर चतुर्भुज है।}
\end{omfigure}

\begin{notation}
जब सिरों से कोई मतलब न हो तब \omterm{def:g10:vectors:vector}{सदिश} को प्रायः एक ही अक्षर से नाम दिया जाता है,
$\vec u$ या $\vec v$। \emph{शून्य \omterm{def:g10:vectors:vector}{सदिश}} $\vec 0 = \vect{AA}$ वह विस्थापन है जो किसी को हिलाता नहीं।
$\vect{AB}$ का \emph{विपरीत} $\vect{BA}$ है: वही लंबाई, उल्टा ढंग। हम $\vect{BA} = -\vect{AB}$ लिखते हैं।
\end{notation}

\section{सदिशों का योग}

\begin{definition}[दो सदिशों का योग]\label{def:g10:vectors:sum}
\emph{योग}\index{सदिश!योग} $\vec u + \vec v$ वह विस्थापन है जो पहले $\vec u$ और फिर $\vec v$
करने से मिलता है।
\end{definition}

\begin{theorem}[शाल का संबंध]\label{thm:g10:vectors:chasles}
किन्हीं तीन बिंदुओं $A$, $B$, $C$ के लिए:
\[
\vect{AB} + \vect{BC} = \vect{AC}.
\]
\end{theorem}

\begin{proof}
$A$ से $B$ तक और फिर $B$ से $C$ तक जाना ऐसा विस्थापन है जो $A$ को $C$
तक ले जाता है: यही विस्थापन $\vect{AC}$ समेटता है।
\end{proof}

\begin{remark}[समांतर चतुर्भुज नियम]
एक ही बिंदु से खींचे गए दो \omterm{def:g10:vectors:vector}{सदिशों} $\vect{AB} + \vect{AC}$ को जोड़ने के लिए समांतर चतुर्भुज $ABDC$
पूरा कीजिए: \omterm{def:g10:vectors:sum}{योग} विकर्ण $\vect{AD}$ है। सचमुच $\vect{AC} = \vect{BD}$, इसलिए शाल के संबंध से $\vect{AB} + \vect{AC} = \vect{AB} + \vect{BD} = \vect{AD}$।
\end{remark}

\begin{omfigure}
\begin{tikzpicture}[scale=1.05]
  \coordinate (A) at (0,0);
  \coordinate (B) at (2.4,0.5);
  \coordinate (C) at (3.4,2.0);
  \draw[-Stealth, omDef, thick] (A) -- (B)
    node[midway, below, font=\small] {$\vec u$};
  \draw[-Stealth, omThm, thick] (B) -- (C)
    node[midway, right, font=\small] {$\vec v$};
  \draw[-Stealth, omMeth, thick] (A) -- (C)
    node[midway, above left, font=\small] {$\vec u + \vec v$};
\end{tikzpicture}\hspace{1.4cm}%
\begin{tikzpicture}[scale=1.05]
  \coordinate (A) at (0,0);
  \coordinate (B) at (2.4,0.5);
  \coordinate (C) at (1.0,1.5);
  \coordinate (D) at (3.4,2.0);
  \draw[omExo, thin] (B) -- (D) -- (C);
  \draw[-Stealth, omDef, thick] (A) -- (B)
    node[midway, below, font=\small] {$\vect{AB}$};
  \draw[-Stealth, omThm, thick] (A) -- (C)
    node[midway, left, font=\small] {$\vect{AC}$};
  \draw[-Stealth, omMeth, thick] (A) -- (D)
    node[pos=0.75, below right, font=\small] {$\vect{AD}$};
  \fill (D) circle (1.2pt) node[right, font=\small] {$D$};
\end{tikzpicture}

{\small \omterm{def:g10:vectors:sum}{योग} के दो चित्र: सिरे-से-पुच्छ (शाल का संबंध, बाएँ) और समांतर
चतुर्भुज नियम (दाएँ)।}
\end{omfigure}

\begin{example}[शाल से सरल करना]\label{ex:g10:vectors:chasles}
$\vect{MN} + \vect{NP} + \vect{PQ}$ को सरल कीजिए: विस्थापनों को जोड़ते जाने पर $\vect{MQ}$ मिलता है। $\vect{AB} - \vect{AC}$ को सरल
कीजिए: अंतर को विपरीत \omterm{def:g10:vectors:vector}{सदिश} के साथ \omterm{def:g10:vectors:sum}{योग} के रूप में फिर से लिखिए,
\[
\vect{AB} - \vect{AC} = \vect{AB} + \vect{CA}
= \vect{CA} + \vect{AB} = \vect{CB}.
\]
\end{example}

\section{सदिश को संख्या से गुणा करना}

\begin{definition}[अदिश गुणज]\label{def:g10:vectors:scalar}
मान लीजिए $\vec u$ एक शून्येतर \omterm{def:g10:vectors:vector}{सदिश} है और $k$ एक \omterm{def:g10:numbers:sets}{वास्तविक संख्या}। \omterm{def:g10:vectors:vector}{सदिश}
$k\vec u$\index{सदिश!अदिश गुणज} की दिशा $\vec u$ जैसी है, लंबाई $\abs{k} \times \norm{\vec u}$ है, और उसका ढंग
$\vec u$ जैसा है यदि $k > 0$ हो, और उल्टा यदि $k < 0$ हो। $k = 0$ के लिए $0\vec u = \vec 0$।
\end{definition}

\begin{definition}[संरेखता]\label{def:g10:vectors:collinear}
दो \omterm{def:g10:vectors:vector}{सदिश} $\vec u$ और $\vec v$ \emph{संरेखी}\index{संरेखी सदिश} तब कहलाते हैं जब उनकी
दिशा एक ही हो, अर्थात् जब किसी वास्तविक $k$ के लिए $\vec v = k \vec u$ हो (या उनमें से एक
$\vec 0$ हो)।
\end{definition}

\begin{remark}
संरेखता समांतरता और संरेखण की सदिश-भाषा है:
\begin{itemize}
  \item रेखाएँ $(AB)$ और $(CD)$ समांतर ठीक तब हैं जब $\vect{AB}$ और $\vect{CD}$ \omterm{def:g10:vectors:collinear}{संरेखी} हों;
  \item बिंदु $A$, $B$, $C$ संरेख ठीक तब हैं जब $\vect{AB}$ और $\vect{AC}$ \omterm{def:g10:vectors:collinear}{संरेखी} हों।
\end{itemize}
\end{remark}

\section{सदिशों के निर्देशांक}

\begin{definition}[सदिश के निर्देशांक]\label{def:g10:vectors:coords}
किसी \omterm{def:g10:coordgeom:system}{निर्देशांक-पद्धति} में \omterm{def:g10:vectors:vector}{सदिश} $\vect{AB}$ के \emph{\omterm{def:g10:coordgeom:system}{निर्देशांक}} वे संख्याएँ हैं
जो विस्थापन का वर्णन करती हैं:
\[
\vect{AB}\,(x_B - x_A,\ y_B - y_A).
\]
\omterm{def:g10:vectors:vector}{सदिश} $\vec u\,(a, b)$ हर बिंदु को क्षैतिज दिशा में $a$ मात्रक और ऊर्ध्वाधर दिशा में
$b$ मात्रक हटा देता है।
\end{definition}

\begin{proposition}[निर्देशांकों के साथ गणना]\label{prop:g10:vectors:coordrules}
मान लीजिए $\vec u\,(a, b)$ और $\vec v\,(c, d)$ हैं, तथा $k \in \R$ है।
\begin{enumerate}
  \item $\vec u = \vec v$ ठीक तब है जब $a = c$ और $b = d$ हों;
  \item $\vec u + \vec v$ के \omterm{def:g10:coordgeom:system}{निर्देशांक} $(a + c,\ b + d)$ हैं;
  \item $k\vec u$ के \omterm{def:g10:coordgeom:system}{निर्देशांक} $(ka,\ kb)$ हैं;
  \item $\norm{\vec u} = \sqrt{a^2 + b^2}$ (\omterm{def:g10:coordgeom:system}{लंबकोणिक पद्धति})।
\end{enumerate}
\end{proposition}

\begin{proof}
1--3 केवल निर्देशांक-दर-निर्देशांक यह दोहराते हैं कि विस्थापन क्या करते हैं:
जैसे पहले $\vec u$ और फिर $\vec v$ करने से कोई बिंदु क्षैतिज दिशा में पहले $a$
और फिर $c$ हटता है, अर्थात् कुल मिलाकर $a + c$। बिंदु 4 \cref{thm:g10:coordgeom:distance} का दूरी सूत्र
है, जो $\vec u$ को दर्शाने वाले किसी रेखाखंड पर लगाया गया है।
\end{proof}

\begin{theorem}[संरेखता की कसौटी]\label{thm:g10:vectors:det}
दो \omterm{def:g10:vectors:vector}{सदिश} $\vec u\,(a, b)$ और $\vec v\,(c, d)$ \omterm{def:g10:vectors:collinear}{संरेखी} हैं यदि और केवल यदि
\[
ad - bc = 0 .
\]
\end{theorem}

\begin{proof}
यदि $\vec v = k\vec u$ हो, तो $c = ka$ और $d = kb$, इसलिए $ad - bc = a(kb) - b(ka) = 0$। यदि $\vec u = k \vec v$ हो, या दोनों में से
कोई \omterm{def:g10:vectors:vector}{सदिश} शून्य हो, तब भी यही बात लागू है।

इसके उलट, मान लीजिए $ad - bc = 0$ है, जहाँ $\vec u \neq \vec 0$ है, मान लीजिए $a \neq 0$ (स्थिति $b \neq 0$
इसी जैसी है)। $k = \frac{c}{a}$ रखिए; तब $c = ka$, और $ad = bc = b(ka)$ से, $a \neq 0$ से भाग देने पर,
$d = kb$ मिलता है। इसलिए $\vec v = k\vec u$।
\end{proof}

\begin{example}\label{ex:g10:vectors:det}
क्या $\vec u\,(3, -2)$ और $\vec v\,(-6, 4)$ \omterm{def:g10:vectors:collinear}{संरेखी} हैं? $ad - bc = 3 \times 4 - (-2)\times(-6) = 12 - 12 = 0$ निकालिए: हाँ, और सचमुच $\vec v = -2\vec u$। $\vec u\,(3, -2)$
और $\vec w\,(1, 5)$ के लिए: $3 \times 5 - (-2) \times 1 = 17 \neq 0$: \omterm{def:g10:vectors:collinear}{संरेखी} नहीं।
\end{example}

\begin{method}[संरेखण और समांतरता]\label{met:g10:vectors:alignment}
यह सिद्ध करने के लिए कि तीन बिंदु $A$, $B$, $C$ संरेख हैं:
\begin{enumerate}
  \item $\vect{AB}$ और $\vect{AC}$ के \omterm{def:g10:coordgeom:system}{निर्देशांक} निकालिए;
  \item \cref{thm:g10:vectors:det} की कसौटी $ad - bc = 0$ जाँचिए;
  \item निष्कर्ष निकालिए: दोनों \omterm{def:g10:vectors:vector}{सदिश} \omterm{def:g10:vectors:collinear}{संरेखी} हैं और उनमें बिंदु $A$
        उभयनिष्ठ है, इसलिए तीनों बिंदु संरेख हैं।
\end{enumerate}
वही गणना $\vect{AB}$ और $\vect{CD}$ के साथ यह सिद्ध करती है कि रेखाएँ $(AB)$ और $(CD)$
समांतर हैं।
\end{method}

\begin{example}\label{ex:g10:vectors:alignment}
मान लीजिए $A(1, 2)$, $B(3, 5)$ और $C(7, 11)$ हैं। तब $\vect{AB}\,(2, 3)$ और $\vect{AC}\,(6, 9)$, और $2 \times 9 - 3 \times 6 = 18 - 18 = 0$: बिंदु
संरेख हैं (वस्तुतः $\vect{AC} = 3\vect{AB}$)।
\end{example}

\begin{omfigure}
\begin{tikzpicture}[scale=0.62]
  \draw[thin, gray!40] (-0.4,-0.4) grid (8.4,5.4);
  \draw[-Stealth] (-0.6,0) -- (8.6,0);
  \draw[-Stealth] (0,-0.6) -- (0,5.6);
  \coordinate (A) at (1,1);
  \coordinate (B) at (3,2);
  \coordinate (C) at (7,4);
  \draw[-Stealth, omDef, very thick] (A) -- (B)
    node[midway, above left=-2pt, font=\small] {$\vect{AB}$};
  \draw[-Stealth, omThm, thick] (A) -- (C)
    node[pos=0.8, below right=-2pt, font=\small] {$\vect{AC} = 3\vect{AB}$};
  \fill (A) circle (2pt) node[above left, font=\small] {$A$};
  \fill (B) circle (2pt) node[above left, font=\small] {$B$};
  \fill (C) circle (2pt) node[above left, font=\small] {$C$};
\end{tikzpicture}

{\small संरेख बिंदु: $\vect{AC}$ $\vect{AB}$ का एक अदिश गुणज है।}
\end{omfigure}

\section{अभ्यास}

\begin{exercise}[$\star$]\label{exo:g10:vectors:1}
शाल के संबंध से सरल कीजिए:
\[
\vect{AB} + \vect{BD}, \qquad
\vect{KL} + \vect{LM} + \vect{MK}, \qquad
\vect{AC} - \vect{BC}, \qquad
\vect{AB} + \vect{CA}.
\]
\end{exercise}

\begin{exercise}[$\star$]\label{exo:g10:vectors:2}
मान लीजिए $A(2, 1)$, $B(5, 3)$, $C(-1, 4)$ हैं। $\vect{AB}$, $\vect{BC}$, $\vect{AC}$ के \omterm{def:g10:coordgeom:system}{निर्देशांक}
निकालिए, और \omterm{def:g10:coordgeom:system}{निर्देशांकों} पर जाँचिए कि $\vect{AB} + \vect{BC} = \vect{AC}$।
\end{exercise}

\begin{exercise}[$\star$]\label{exo:g10:vectors:3}
मान लीजिए $\vec u\,(2, -3)$ और $\vec v\,(-1, 4)$ हैं। $\vec u + \vec v$, $3\vec u$, $2\vec u - \vec v$ के \omterm{def:g10:coordgeom:system}{निर्देशांक} बताइए, और
$\norm{\vec u}$ निकालिए।
\end{exercise}

\begin{exercise}[$\star$]\label{exo:g10:vectors:4}
हर स्थिति में बताइए कि $\vec u$ और $\vec v$ \omterm{def:g10:vectors:collinear}{संरेखी} हैं या नहीं:
\[
\text{(a) } \vec u\,(4, 6),\ \vec v\,(6, 9);
\qquad
\text{(b) } \vec u\,(2, -5),\ \vec v\,(-4, 10);
\qquad
\text{(c) } \vec u\,(3, 1),\ \vec v\,(1, 3).
\]
\end{exercise}

\begin{exercise}[$\star$]\label{exo:g10:vectors:5}
मान लीजिए $A(0, 2)$, $B(4, 0)$, $C(6, 3)$ हैं। वह बिंदु $D$ ज्ञात कीजिए जिसके लिए
$ABCD$ समांतर चतुर्भुज हो, अर्थात् $\vect{AB} = \vect{DC}$ हो।
\end{exercise}

\begin{exercise}[$\star\star$]\label{exo:g10:vectors:6}
मान लीजिए $A(-2, 1)$, $B(1, 3)$, $C(4, -1)$ हैं।
\begin{enumerate}
  \item वह बिंदु $M$ ज्ञात कीजिए जिसके लिए $\vect{AM} = \vect{AB} + \vect{AC}$ हो।
  \item चतुर्भुज $ABMC$ कौन-सा है? पुष्ट कीजिए।
\end{enumerate}
\end{exercise}

\begin{exercise}[$\star\star$]\label{exo:g10:vectors:7}
मान लीजिए $A(1, -1)$, $B(4, 1)$, $C(10, 5)$ हैं। क्या बिंदु $A$, $B$, $C$ संरेख
हैं? यही प्रश्न $A(0, 3)$, $B(2, 2)$, $C(8, -1)$ के लिए।
\end{exercise}

\begin{exercise}[$\star\star$]\label{exo:g10:vectors:8}
मान लीजिए $A(1, 2)$, $B(5, 4)$, $C(6, 1)$, $D(2, -1)$ हैं। दिखाइए कि $(AB)$ और $(CD)$ समांतर
हैं, फिर यह कि $ABCD$ समांतर चतुर्भुज है। कौन-सी जाँच दूसरी को अपने आप सिद्ध
कर देती है?
\end{exercise}

\begin{exercise}[$\star\star$]\label{exo:g10:vectors:9}
मान लीजिए $ABC$ एक त्रिभुज है, और $I$ वह बिंदु है जो $\vect{AI} = \frac23 \vect{AB}$ से परिभाषित
होता है। $A(0,0)$, $B(6, 3)$, $C(2, 5)$ लेकर (ताकि गणनाएँ सरल रहें):
\begin{enumerate}
  \item $I$ के \omterm{def:g10:coordgeom:system}{निर्देशांक} निकालिए;
  \item मान लीजिए $J$ ऐसा है कि $\vect{CJ} = \frac23\vect{CB}$; $J$ के \omterm{def:g10:coordgeom:system}{निर्देशांक} निकालिए;
  \item दिखाइए कि $(IJ)$ $(AC)$ के समांतर है।
\end{enumerate}
\end{exercise}

\begin{exercise}[$\star\star\star$]\label{exo:g10:vectors:10}
मान लीजिए $ABCD$ कोई चतुर्भुज है, और $I$, $J$, $K$, $L$ क्रमशः
$[AB]$, $[BC]$, $[CD]$, $[DA]$ के \omterm{prop:g10:coordgeom:midpoint}{मध्यबिंदु} हैं। \omterm{def:g10:coordgeom:system}{निर्देशांकों} $A(x_A, y_A)$ आदि का
उपयोग करके दिखाइए कि $\vect{IJ} = \vect{LK}$, और निष्कर्ष निकालिए कि \emph{किसी भी} चतुर्भुज
की भुजाओं के \omterm{prop:g10:coordgeom:midpoint}{मध्यबिंदु} एक समांतर चतुर्भुज बनाते हैं।
\end{exercise}

\section{समस्या: तीरों का बीजगणित}

\begin{problem}[{सप्ताहांत समस्या --- शाल की कसरत, केन्द्रक की दो पंक्तियों
में नई उपपत्ति, और साम्यावस्था में बल: सदिश असल में किस काम आते हैं}]
\label{pb:g10:vectors:1}
कक्षा 8 ने सिद्ध किया था कि त्रिभुज की माध्यिकाएँ अपनी लंबाई के दो-तिहाई पर
मिलती हैं --- और उसके लिए त्रिभुज के भीतर एक चतुर समांतर चतुर्भुज छिपाया गया
था (माध्यमिक विद्यालय खंड की केन्द्रक वाली सप्ताहांत समस्या)। \omterm{def:g10:vectors:vector}{सदिश} उसे दो
पंक्तियों में फिर सिद्ध कर देते हैं, और संगामिता मुफ़्त में दे देते हैं। यही
इस नए बीजगणित का काम है: प्रमेय तीरों के साथ की गई गणनाएँ बन जाती हैं। यह
समस्या पहले वह गणना सधाती है (सबसे बढ़कर शाल की), फिर दो पंक्ति वाली उपपत्ति
देती है, फिर वारिन्यों (\cref{exo:g10:vectors:10}) को आगे बढ़ाती है, और वहाँ समाप्त होती है
जहाँ \omterm{def:g10:vectors:vector}{सदिशों} का जन्म हुआ था: बल और वेग।

\medskip
\noindent\textbf{भाग I --- शाल की कसरत।}
\begin{enumerate}
  \item शाल के संबंध (\cref{thm:g10:vectors:chasles}) से सरल कीजिए:
        $\vect{AB} + \vect{BC} + \vect{CD}$; \quad
        $\vect{MA} - \vect{MB}$; \quad
        $\vect{AB} + \vect{CD} + \vect{BC} + \vect{DA}$।
  \item मूल बिंदु वाली तरकीब: दिखाइए कि \emph{किसी भी} बिंदु $O$ के लिए
        $\vect{AB} = \vect{OB} - \vect{OA}$।
  \item $[AB]$ के \omterm{prop:g10:coordgeom:midpoint}{मध्यबिंदु} $I$ के दोनों अभिलक्षण सिद्ध कीजिए:
        \[
        \vect{IA} + \vect{IB} = \vec 0
        \qquad\text{और}\qquad
        \vect{OI} = \tfrac12\left(\vect{OA} + \vect{OB}\right)
        \ \text{किसी भी } O \text{ के लिए} .
        \]
  \item $\vect{AB} = \vect{DC}$ का उपयोग करके $A(-1, 2)$, $B(3, 4)$, $C(6, 0)$ वाले समांतर चतुर्भुज $ABCD$
        का चौथा शीर्ष $D$ फिर से ज्ञात कीजिए --- और \cref{pb:g10:coordgeom:1} के \omterm{prop:g10:coordgeom:midpoint}{मध्यबिंदु}
        वाली तरकीब से मिले उत्तर से तुलना कीजिए।
  \item समतल के हर बिंदु में एक नियत \omterm{def:g10:vectors:vector}{सदिश} $\vec u$ जोड़ने से कौन-सा रूपांतरण
        होता है --- वही जो माध्यमिक विद्यालय खंड की दो-दर्पण वाली सप्ताहांत
        समस्या में दो दर्पणों से और अर्ध-घूर्णन वाली सप्ताहांत समस्या में दो
        अर्ध-घूर्णनों से खोजा गया था? $\vec u = (3, -1)$ के अधीन $(1, 2)$ का \omterm{def:g10:functions:function}{प्रतिबिंब}
        निकालिए।
\end{enumerate}

\medskip
\noindent\textbf{भाग II --- केन्द्रक, तीसरी उपपत्ति।}
त्रिभुज $ABC$ के बिंदु $G$ को इस सुंदर शर्त से परिभाषित कीजिए:
\[
\vect{GA} + \vect{GB} + \vect{GC} = \vec 0 .
\]
\begin{enumerate}[resume]
  \item मूल बिंदु वाली तरकीब से दिखाइए कि यह शर्त $G$ को अद्वितीय रूप से
        तय कर देती है: किसी भी $O$ के लिए $\vect{OG} = \frac13\left(\vect{OA} + \vect{OB} +
        \vect{OC}\right)$ --- अर्थात् $G$ के
        \omterm{def:g10:coordgeom:system}{निर्देशांक} शीर्षों के \omterm{def:g10:coordgeom:system}{निर्देशांकों} के \emph{माध्य} हैं (माध्यमिक
        विद्यालय खंड की केन्द्रक वाली सप्ताहांत समस्या ने इसे संख्याओं में
        देखा था)।
  \item $\vect{AG} = \frac13\left(\vect{AB} +
        \vect{AC}\right)$ सिद्ध कीजिए, और $[BC]$ के \omterm{prop:g10:coordgeom:midpoint}{मध्यबिंदु} को $K$ लिखकर $\vect{AK} = \frac12\left(\vect{AB} +
        \vect{AC}\right)$।
        निष्कर्ष निकालिए: $\vect{AG} = \frac23\,\vect{AK}$ --- दो-तिहाई वाली प्रमेय, फिर से सिद्ध।
  \item $B$ से और $C$ से खींची गई माध्यिकाएँ बिना किसी और गणना के
        \emph{उसी} $G$ से क्यों गुज़रती हैं? इस प्रमेय की तीनों उपलब्ध
        उपपत्तियों --- कक्षा 8 का छिपा समांतर चतुर्भुज, \omterm{def:g10:coordgeom:system}{निर्देशांक}, और
        \omterm{def:g10:vectors:vector}{सदिश} --- की तुलना एक-एक वाक्य में कीजिए।
  \item संख्यात्मक जाँच: $A(1, 1)$, $B(5, 2)$, $C(3, 6)$। $G$ निकालिए, फिर $K$,
        फिर संरेखता की कसौटी (\cref{thm:g10:vectors:det}) से $\vect{AG} = \frac23 \vect{AK}$ सत्यापित कीजिए।
  \item भौतिकी $G$ को शीर्षों पर रखे तीन बराबर द्रव्यमानों के संतुलन-बिंदु
        के रूप में पढ़ती है। $A$ का द्रव्यमान दुगुना कर दीजिए (द्रव्यमान
        $2, 1, 1$): तब संतुलन-बिंदु $\vect{OG'} = \frac{2\vect{OA} + \vect{OB} +
        \vect{OC}}{4}$ हो जाता है। प्रश्न 9 वाले त्रिभुज के
        लिए $G'$ निकालिए और उसकी स्थिति की $G$ से तुलना कीजिए।
\end{enumerate}

\medskip
\noindent\textbf{भाग III --- संरेखता काम पर।}
\begin{enumerate}[resume]
  \item क्या $\vec u(3, -2)$ और $\vec v(-6, 4)$ \omterm{def:g10:vectors:collinear}{संरेखी} हैं? और $\vec w(2, 5)$, $\vec z(4, 9)$? सारणिक कसौटी से
        तय कीजिए।
  \item क्या बिंदु $A(1, 2)$, $B(3, 5)$, $C(7, 11)$ संरेख हैं?
  \item वारिन्यों से आगे (\cref{exo:g10:vectors:10}): किसी भी चतुर्भुज $ABCD$ में
        \emph{द्वि-माध्यिकाएँ} सामने-सामने की भुजाओं के मध्यबिंदुओं को
        जोड़ती हैं ($[IK]$ और $[JL]$)। स्थिति \omterm{def:g10:vectors:vector}{सदिशों} ($\vect{OI} = \frac12(\vect{OA} +
        \vect{OB})$ आदि) का उपयोग
        करके दिखाइए कि दोनों द्वि-माध्यिकाओं का \omterm{prop:g10:coordgeom:midpoint}{मध्यबिंदु} एक ही है, जिसका
        स्थिति \omterm{def:g10:vectors:vector}{सदिश} $\frac14\left(\vect{OA} + \vect{OB} + \vect{OC} +
        \vect{OD}\right)$ है: यही चतुर्भुज का केन्द्रक है।
  \item थेल्स एक पंक्ति में: $M$ और $N$ वे बिंदु हैं जो $\vect{AM} = \frac13\vect{AB}$ और $\vect{AN} = \frac13\vect{AC}$
        से परिभाषित हैं। $\vect{MN}$ को $\vect{BC}$ के पदों में निकालिए और समांतरता
        तथा अनुपात का निष्कर्ष निकालिए।
  \item प्रश्न 14 को किसी भी अनुपात $k$ तक सामान्य कीजिए, और एक वाक्य में
        बताइए कि सदिश-कलन \omterm{prop:g10:coordgeom:midpoint}{मध्यबिंदु} प्रमेय और थेल्स प्रमेय दोनों को एक साथ
        कैसे अपने भीतर समेट लेता है।
\end{enumerate}

\medskip
\noindent\textbf{भाग IV --- बल और वेग।}
\begin{enumerate}[resume]
  \item एक नाव को $3$ km/h की धारा ठीक पूर्व की ओर धकेल रही है, जबकि
        उसमें $4$ km/h से ठीक उत्तर की ओर चप्पू चलाया जा रहा है। परिणामी
        वेग \omterm{def:g10:vectors:vector}{सदिश}, उसका परिमाण और नाव का वास्तविक मार्ग बताइए। (एक बहुत
        पुराना त्रिक प्रकट होता है।)
  \item एक छल्ले पर तीन बल लगते हैं: $\vec u(2, 1)$, $\vec v(-3, 2)$, $\vec w(1, -3)$। परिणामी निकालिए।
        छल्ला क्या करता है? और किसी भी \emph{बंद} बहुभुज के अनुदिश तीरों का
        \omterm{def:g10:vectors:sum}{योग} सदा $\vec 0$ क्यों होता है (शाल)?
  \item एक हुक पर दो बल $\vec u(4, -1)$ और $\vec v(-1, 3)$ लगते हैं। कौन-सा तीसरा बल हुक को
        साम्यावस्था में रखेगा?
  \item एक तैराक नदी को सीधे पार करने की दिशा में $2$ m/s से तैरता है;
        धारा $1.5$ m/s से बहती है। परिणामी वेग और उसका परिमाण बताइए। नदी
        $50$ m चौड़ी है: पार करने में कितना समय लगता है, और तैराक धारा की
        दिशा में कितनी दूर जा लगता है?
  \item समापन: एक ही वस्तु, तीन वेश --- विस्थापन का तीर, निर्देशांक-युग्म,
        और बल या वेग। इस समस्या ने सदिश-बीजगणित से जिन तीन चिरपरिचित
        प्रमेयों को फिर सिद्ध किया, उन्हें गिनाइए, और उस एक ज्यामितीय राशि
        का नाम बताइए जिसे इस अध्याय के \omterm{def:g10:vectors:vector}{सदिश} अब भी नाप नहीं सकते --- उसे
        साधने वाला औज़ार कक्षा 11 में अदिश गुणनफल के साथ आता है।
\end{enumerate}
\end{problem}
